\documentclass[11pt]{article}
\usepackage[margin=1in]{geometry}
\usepackage{amsmath,amssymb}
\usepackage{xcolor}
\usepackage{booktabs}
\usepackage{enumitem}
\usepackage{fancyhdr}
\usepackage{tcolorbox}
\tcbuselibrary{skins,breakable}
\usepackage{hyperref}
\hypersetup{colorlinks=true, linkcolor=blue, urlcolor=blue}

\definecolor{primary}{RGB}{30,64,140}
\definecolor{tipbg}{RGB}{235,244,255}
\definecolor{examplebg}{RGB}{240,248,235}
\definecolor{formulabg}{RGB}{255,247,224}
\definecolor{warnbg}{RGB}{255,236,236}

\pagestyle{fancy}
\fancyhf{}
\renewcommand{\headrulewidth}{0.4pt}
\cfoot{\thepage}

\newtcolorbox{tipbox}[1][Tip]{
  colback=tipbg, colframe=primary, fonttitle=\bfseries,
  title=#1, breakable, rounded corners
}
\newtcolorbox{examplebox}[1][Example]{
  colback=examplebg, colframe=primary!70!black, fonttitle=\bfseries,
  title=#1, breakable, rounded corners
}
\newtcolorbox{formulabox}[1][Formula]{
  colback=formulabg, colframe=primary!70!black, fonttitle=\bfseries,
  title=#1, breakable, rounded corners
}
\newtcolorbox{warnbox}[1][Warning]{
  colback=warnbg, colframe=red!60!black, fonttitle=\bfseries,
  title=#1, breakable, rounded corners
}


\title{\color{primary}\textbf{Time \& Work --- Practice Questions}\\[4pt]\large Medium to Difficult, With Full Solutions}
\author{Aptitude Grind}
\date{}

\lhead{\small Time \& Work --- Practice Questions}
\rhead{\small Aptitude Grind}

\begin{document}
\maketitle
\thispagestyle{fancy}

\noindent\textit{Difficulty is marked as \textbf{[M]} Medium or \textbf{[H]} Hard. Solve on your own
first, then check the solution. Solutions use the methods from \texttt{01-fundamentals.tex},
\texttt{02-fast-tricks-and-shortcuts.tex}, and \texttt{03-question-pattern-catalog.tex}. Every
numeric answer below was cross-checked independently using exact fraction arithmetic (Python's
\texttt{fractions.Fraction}) before being finalized.}

\bigskip
\hrule
\bigskip

\subsection*{Q1 [M] --- Two People Working Together (Pattern A1)}
A can complete a job alone in 12 days, B alone in 24 days. In how many days can they finish it
working together?

\begin{tipbox}[Solution]
\[
\text{Together} = \frac{12 \times 24}{12+24} = \frac{288}{36} = 8
\]
\textbf{Answer: 8 days}
\end{tipbox}

\subsection*{Q2 [M] --- Efficiency Ratio (Pattern B1)}
The efficiency ratio of A to B is 3:5. A alone takes 20 days to complete a job. Find the number of
days B alone would take, and the number of days they would take working together.

\begin{tipbox}[Solution]
Efficiency ratio 3:5 means A is the \emph{less} efficient one, so A (efficiency 3) takes more time
than B (efficiency 5). Since time is inversely proportional to efficiency:
\begin{align*}
\text{Time ratio } A:B &= 5:3 \quad \text{(reverse of efficiency ratio)} \\
A = 20 \text{ days} \;\Rightarrow\; B &= 20 \times 3/5 = 12 \text{ days} \\
\text{Together} &= \frac{20 \times 12}{20+12} = \frac{240}{32} = 7.5 \text{ days}
\end{align*}
\textbf{Answer: B alone = 12 days; together = 7.5 days}
\end{tipbox}

\subsection*{Q3 [M] --- The Difference Formula (Pattern \S5 of \texttt{01-fundamentals.tex})}
A alone can complete a job in 30 days. A and B together can complete it in 12 days. In how many days
can B alone complete the job?

\begin{tipbox}[Solution]
\[
\frac{1}{B} = \frac{1}{12} - \frac{1}{30} = \frac{5}{60} - \frac{2}{60} = \frac{3}{60} = \frac{1}{20}
\]
\textbf{Answer: B alone = 20 days}
\end{tipbox}

\subsection*{Q4 [M] --- LCM Method, Three Workers (Pattern A2)}
A, B, and C can individually complete a job in 12, 15, and 20 days respectively. In how many days
will they finish it working together?

\begin{tipbox}[Solution]
\begin{align*}
\text{LCM}(12,15,20) &= 60 \text{ units of total work} \\
A\text{'s rate} = 60/12 = 5,\quad B\text{'s rate} &= 60/15 = 4,\quad C\text{'s rate} = 60/20 = 3
\text{ (units/day)} \\
\text{Combined rate} &= 5+4+3 = 12 \text{ units/day} \\
\text{Time} &= 60/12 = 5 \text{ days}
\end{align*}
\textbf{Answer: 5 days}
\end{tipbox}

\subsection*{Q5 [M] --- Pipes Filling and Emptying Together (Pattern E1)}
A pipe can fill a tank in 6 hours; another pipe can empty the full tank in 9 hours. If both pipes are
opened together, in how much time will the tank be filled?

\begin{tipbox}[Solution]
\[
\text{Net rate} = \frac{1}{6} - \frac{1}{9} = \frac{3}{18} - \frac{2}{18} = \frac{1}{18}
\Rightarrow \text{Time} = 18 \text{ hours}
\]
\textbf{Answer: 18 hours}
\end{tipbox}

\subsection*{Q6 [H] --- Pipe Closed Partway Through (Pattern E2)}
Fill pipe A can fill a tank in 10 hours; outlet pipe B can empty it in 15 hours. Both are opened
together, but B is closed after 3 hours. Find the total time to fill the tank.

\begin{tipbox}[Solution]
\begin{align*}
\text{Net rate (both open)} &= 1/10 - 1/15 = 1/30 \\
\text{Work done in first 3 hours} &= 3 \times 1/30 = 1/10 \\
\text{Remaining work} &= 1 - 1/10 = 9/10 \\
\text{A alone (rate } 1/10\text{/hr) needs:} & (9/10) \div (1/10) = 9 \text{ more hours} \\
\text{Total time} &= 3 + 9 = 12 \text{ hours}
\end{align*}
\textbf{Answer: 12 hours}
\end{tipbox}

\begin{warnbox}[Why This Is a Common Trap]
People sometimes keep using the \emph{net} rate for the whole problem, forgetting that once the
outlet closes, only the inlet's rate applies for the rest of the time. Always split into phases the
moment any pipe's status changes.
\end{warnbox}

\subsection*{Q7 [H] --- Worker Joining Midway (Pattern C1)}
A can complete a job alone in 20 days. A works alone for 5 days, then B (who alone would take 30
days) joins. How many total days does the job take to finish?

\begin{tipbox}[Solution]
\begin{align*}
\text{Work done by A in 5 days} &= 5/20 = 1/4 \\
\text{Remaining work} &= 3/4 \\
\text{Combined rate (A + B)} &= 1/20 + 1/30 = 3/60 + 2/60 = 1/12 \\
\text{Extra days needed} &= (3/4) \div (1/12) = 9 \\
\text{Total days} &= 5 + 9 = 14
\end{align*}
\textbf{Answer: 14 days}
\end{tipbox}

\subsection*{Q8 [H] --- Workers Leaving Midway (Pattern C2)}
36 workers can complete a job in 20 days. After 5 days, 6 workers leave the job. In how many more
days will the remaining workers finish it?

\begin{tipbox}[Solution]
\begin{align*}
\text{Total work} &= 36 \times 20 = 720 \text{ worker-days} \\
\text{Work done in 5 days} &= 36 \times 5 = 180 \text{ worker-days} \\
\text{Remaining work} &= 720 - 180 = 540 \text{ worker-days} \\
\text{Workers remaining} &= 36 - 6 = 30 \\
\text{Days needed} &= 540/30 = 18
\end{align*}
\textbf{Answer: 18 more days}
\end{tipbox}

\subsection*{Q9 [M] --- Alternate-Day Working (Pattern F1)}
A can complete a job alone in 6 days, B alone in 9 days. Working on alternate days, starting with A,
in how many days will the job be finished?

\begin{tipbox}[Solution]
\[
\text{2-day cycle work} = \frac{1}{6} + \frac{1}{9} = \frac{3}{18} + \frac{2}{18} = \frac{5}{18}
\]
Three full cycles (6 days) complete $3 \times 5/18 = 15/18 = 5/6$ of the work. Remaining work $=1/6$.
Day 7 is A's turn again: A's rate is $1/6$, so A finishes exactly the remaining $1/6$ in \textbf{one
more full day}.
\[
\text{Total} = 6 \text{ (three full cycles)} + 1 \text{ (A's final day)} = 7 \text{ days}
\]
\textbf{Answer: 7 days}
\end{tipbox}

\subsection*{Q10 [H] --- Men/Women Equivalence (Pattern D1)}
3 men do the same amount of work as 5 women (equal efficiency comparison). 3 men alone can complete a
job in 20 days. In how many days can 6 men and 10 women together complete the same job?

\begin{tipbox}[Solution]
\begin{align*}
\text{3 men alone: 20 days} &\Rightarrow 1 \text{ man's rate} = 1/(3\times20) = 1/60 \\
3 \text{ men} = 5 \text{ women (equal output)} &\Rightarrow 1 \text{ woman's rate} =
(3\times1/60)/5 = 1/100 \\
6 \text{ men} + 10 \text{ women rate} &= 6\times(1/60) + 10\times(1/100) = 1/10 + 1/10 = 1/5 \\
\text{Time} &= 5 \text{ days}
\end{align*}
\textbf{Answer: 5 days}
\end{tipbox}

\begin{tipbox}[Note]
6 men and 10 women is exactly \emph{double} the ``3 men = 5 women'' equivalence group, so it makes
sense the answer is a clean number --- doubling the workforce on a fixed job always halves the time,
and here the group ratio was scaled up perfectly.
\end{tipbox}

\subsection*{Q11 [M] --- Wages Based on Work (Pattern G1)}
A can complete a job alone in 8 days, B alone in 12 days. They work together and are jointly paid
Rs.~2500. Find each person's share.

\begin{tipbox}[Solution]
\begin{align*}
\text{Rate ratio } A:B &= 1/8 : 1/12 = 3:2 \quad \text{(multiply both by 24)} \\
A\text{'s share} &= 2500 \times 3/5 = Rs.~1500 \\
B\text{'s share} &= 2500 \times 2/5 = Rs.~1000
\end{align*}
\textbf{Answer: A = Rs.~1500, B = Rs.~1000}
\end{tipbox}

\subsection*{Q12 [M] --- Efficiency Decrease $\to$ Time Change (Pattern J2)}
A's efficiency decreases by 25\%. If A originally took 18 days to complete a job at full efficiency,
how many days will A now take?

\begin{tipbox}[Solution]
\[
\text{New Time} = \text{Old Time} \times \frac{100}{100-25} = 18 \times \frac{100}{75} = 24
\]
\textbf{Answer: 24 days}
\end{tipbox}

\subsection*{Q13 [H] --- Variable Efficiency Mid-Job (Pattern J1)}
Working at a constant rate, A can finish a job in 12 days. A works at the normal rate for the first 3
days, then works at \textbf{triple} efficiency for the rest of the job. How many total days does the
job take?

\begin{tipbox}[Solution]
\begin{align*}
\text{Normal rate} &= 1/12 \\
\text{Work done in first 3 days} &= 3 \times 1/12 = 1/4 \\
\text{Remaining work} &= 3/4 \\
\text{Tripled rate} &= 3 \times 1/12 = 1/4 \\
\text{Extra days needed} &= (3/4) \div (1/4) = 3 \\
\text{Total days} &= 3 + 3 = 6
\end{align*}
\textbf{Answer: 6 days}
\end{tipbox}

\subsection*{Q14 [H] --- Percentage of Work Completed (Pattern I1)}
A can complete 90\% of a job in 18 days. In how many days can A alone complete the entire job? B then
finishes the remaining 10\% of the job in 3 days --- in how many days could B alone complete the
entire job?

\begin{tipbox}[Solution]
\begin{align*}
A: 90\% \text{ in 18 days} &\Rightarrow 100\% \text{ in } 18/0.9 = 20 \text{ days} \\
B: 10\% \text{ in 3 days} &\Rightarrow 100\% \text{ in } 3/0.1 = 30 \text{ days}
\end{align*}
\textbf{Answer: A alone = 20 days; B alone = 30 days}
\end{tipbox}

\begin{warnbox}[Why This Is a Common Trap]
People sometimes try to combine A and B's \emph{individual} full-job times directly (e.g., averaging
20 and 30) to answer a different question about their combined time. Here the question only asks for
each person's \textbf{individual} full-job time --- scale each person's given fraction
independently; don't mix the two workers' numbers together unless the question specifically asks for
a combined rate.
\end{warnbox}

\section*{Answer Key Summary}
\begin{formulabox}[Answers]
\centering
\begin{tabular}{ll}
\toprule
\textbf{Question} & \textbf{Answer} \\
\midrule
Q1 & 8 days \\
Q2 & B = 12 days; together = 7.5 days \\
Q3 & 20 days \\
Q4 & 5 days \\
Q5 & 18 hours \\
Q6 & 12 hours \\
Q7 & 14 days \\
Q8 & 18 more days \\
Q9 & 7 days \\
Q10 & 5 days \\
Q11 & A = Rs.~1500, B = Rs.~1000 \\
Q12 & 24 days \\
Q13 & 6 days \\
Q14 & A = 20 days, B = 30 days \\
\bottomrule
\end{tabular}
\end{formulabox}

\end{document}
